\section{\texorpdfstring{映射横截性与 $\bar\partial$ 的线性化}{映射横截性与柯西--黎曼算子的线性化}}
\label{sec:mapping-transversality-linearization}
\stfhSubIndexRangeStart{I292}{heng^jie^xing}{横截性（transversality）}{ying^she^heng^jie^xing}{映射横截性（mapping transversality）}
\setcounter{equation}{6}

现在我们希望证明 $\widetilde{\mathcal M}(J;\beta)$ 具有光滑流形结构。我们希望通过适当选取殆复结构
$J\in\mathcal J_\omega$ 来获得这种\emph{横截性}。

为此，我们把 $J$ 看作参数，把 $\mathcal J_\omega$ 看作参数空间；这里
$\mathcal J_\omega$ 是与 $\omega$ 相容的殆复结构空间。在本节余下部分，我们固定 $j$，从而略去对
$j$ 的依赖。由于赋予 $C^\infty$ 拓扑的 $\mathcal J_\omega$ 不是巴拿赫流形，在技术上还需考虑它的
$C^\ell$ 类似物，记作 $\mathcal J_\omega^\ell$，其中 $\ell\geq1$。于是
$\mathcal J_\omega^\ell$ 具有巴拿赫流形结构。

给定 $(k,p)$、$\ell$ 及类 $\beta\in\pi_2(M)$，考虑丛
\[
 \mathcal H''(\beta)\longrightarrow
 \mathcal F^{k,p}(\beta)\times\mathcal J_\omega^\ell,
\]
其在 $(u,J)$ 处的纤维为 $\mathcal H_{(u,J)}^{(0,1)}$，即
\[
 \mathcal H''(\beta)=
 \bigcup_{(u,J)\in\mathcal F^{k,p}(\beta)\times\mathcal J_\omega^\ell}
 \{(u,J)\}\times\mathcal H_{(u,J)}^{(0,1)}.
\]
把映射
\[
 \bar\partial:\mathcal F^{k,p}(\beta)\times\mathcal J_\omega^\ell
 \longrightarrow\mathcal H''(\beta),
 \qquad \bar\partial(u,J)=\bar\partial_Ju,
\]
看作这个丛的截面。我们要研究它在每个
$(u,J)\in\bar\partial^{-1}(0)$ 处的线性化。记零点集为
\[
 \widetilde{\mathcal M}(\beta)=
 \{(u,J)\in\mathcal F^{k,p}(\beta)\times\mathcal J_\omega^\ell
 \mid \bar\partial_{(j,J)}u=0\},
\]
称为普适模空间。以
$\pi_2:\mathcal F^{k,p}(\beta)\times\mathcal J_\omega^\ell
\to\mathcal J_\omega^\ell$ 表示投影，则
\[
 \widetilde{\mathcal M}(J;\beta)=\pi_2^{-1}(J)\cap
 \widetilde{\mathcal M}(\beta).
\]

下面的定理概括了伪全纯曲线模空间研究的主要方案。

\begin{theorem}\label{thm:mapping-transversality-main}
设 $0<\ell<k-2/p$，并以
$\widetilde{\mathcal M}^{\mathrm{inj}}(\beta)$ 表示子集
\[
 \widetilde{\mathcal M}^{\mathrm{inj}}(\beta)=
 \{(u,J)\in\widetilde{\mathcal M}(\beta)
 \mid u\text{ 在某处单射}\}.
\]
则：
\begin{enumerate}[label=(\arabic*)]
 \item $\widetilde{\mathcal M}^{\mathrm{inj}}(\beta)$ 是无穷维 $C^\ell$ 巴拿赫流形；
 \item 投影
 \[
  \Pi_\beta:=\left.\pi_2\right|_{\widetilde{\mathcal M}^{\mathrm{inj}}(\beta)}:
  \widetilde{\mathcal M}^{\mathrm{inj}}(\beta)\longrightarrow\mathcal J_\omega^\ell
 \]
 是弗雷德霍姆映射，且其指标与某个（因而任意）
 $u\in\widetilde{\mathcal M}^{\mathrm{inj}}(\beta)$ 处
 $D\bar\partial_J(u)$ 的指标相同。
\end{enumerate}
\end{theorem}

本节余下部分用于证明这个定理。对第一项，只需证明在 $u$ 为在某处单射曲线的每个 $(u,J)$ 处，协变线性化
\[
 D\bar\partial(u,J):T_u\mathcal F^{k,p}(\beta)\oplus
 T_J\mathcal J_\omega^\ell\longrightarrow\mathcal H''_{(u,J)}
\]
是满射。为避免记号 $\bar\partial$ 引起混淆，在以下计算中记
$\Phi=\bar\partial$。注意
\[
 D\Phi(u,J)=D_1\Phi(u,J)+D_2\Phi(u,J),
\]
其中相伴的偏导数为
\[
 D_1\Phi(u,J):T_u\mathcal F^{k,p}(\beta)\to\mathcal H''_{(u,J)},
 \qquad
 D_2\Phi(u,J):T_J\mathcal J_\omega^\ell\to\mathcal H''_{(u,J)}.
\]
直接计算得到
\begin{equation}
 D_2\Phi(u,J)(B)=\frac12 B\circ du\circ j.
 \label{eq:variation-almost-complex-structure}
\end{equation}
显然按定义有 $D_1\Phi(u,J)=D_u\bar\partial_{(j,J)}$。我们已经计算过
$D_u\bar\partial_{(j,J)}$，并证明它是形式椭圆的。因此它向 $W^{k,p}$ 的延拓（仍记作
$D_1\Phi(u,J)$）是弗雷德霍姆的；特别地，$\operatorname{coker}D_1\Phi(u,J)$ 有限维。
又因为 $\mathcal H''_{(u,J)}$ 包含 $\operatorname{Im}D_1\Phi(u,J)$，这说明
$\operatorname{Im}D\Phi(u,J)$ 本身是 $\mathcal H''_{(u,J)}$ 中有限余维的闭子空间。

所以，为证明 $D\Phi(u,J)$ 满射，只需证明
$\operatorname{coker}D\Phi(u,J)=\{0\}$。哈恩--巴拿赫定理把问题化为：对任意连续线性泛函
$\eta:\mathcal H''_{(u,J)}\to\mathbb R$，若
$\eta|_{\operatorname{Im}D\Phi(u,J)}=0$，则 $\eta=0$。这将由下面更强的引理推出。

\begin{lemma}\label{lem:annihilator-universal-linearization}
若 $L^p(\Lambda_{(j,J)}^1(u^*TM))$ 上的任意连续线性泛函 $\eta$ 在
$D\Phi(u,J)(C^\infty(\Gamma(u^*TM)))$ 上消失，则 $\eta=0$。
\end{lemma}

暂且推迟证明这个引理，先完成满射性的证明。令 $k,p$ 为上述任意常数，并通过连续嵌入
\[
 \mathcal H''_{(u,J)}=W^{k-1,p}(\Lambda_{(j,J)}^1(u^*TM))
 \hookrightarrow L^p(\Lambda_{(j,J)}^1(u^*TM))
\]
把 $\mathcal H''_{(u,J)}$ 看作 $L^p$ 的子空间。设 $\eta$ 是
$W^{k-1,p}(\Lambda_{(j,J)}^1(u^*TM))$ 上在
$\operatorname{Im}D\Phi(u,J)$ 上消失的连续线性泛函。由于
\[
 \operatorname{Im}D\Phi(u,J)\supset
 D\Phi(u,J)(C^\infty(\Gamma(u^*TM))),
\]
引理 10.4.2 给出 $\eta$ 在 $L^p$ 上为零，因而在 $W^{k-1,p}$ 上也为零。除引理 10.4.2 的证明外，这就完成了满射性的证明。

本节下面将证明引理 10.4.2。

\begin{proof}[引理 10.4.2 的证明]
利用 $\Sigma$ 上的积分配对，我们作如下等同：
\[
 \bigl(L^p(\Lambda^{(0,1)}(u^*TM))\bigr)^*
 \cong L^q(\Lambda^{(1,0)}(u^*TM)),
\]
其中 $1/p+1/q=1$ 且 $1<q<2$；配对
\[
 L^p(\Lambda^{(0,1)}(u^*TM))\times
 L^q(\Lambda^{(1,0)}(u^*TM))\longrightarrow\mathbb R,
 \qquad (f,g)\longmapsto\int_\Sigma\langle f,g\rangle
\]
是非退化的。我们要证明：若
$\eta\in L^q(\Lambda_{(j,J)}^{(1,0)}(u^*TM))$ 满足
\begin{equation}
 \int_\Sigma\langle D\Phi(u,J)(\xi,B),\eta\rangle=0
 \label{eq:annihilator-full-linearization}
\end{equation}
对所有 $\xi\in C^\infty(u^*TM)$ 及 $B\in C^\ell$ 成立，则 $\eta\equiv0$。

方程\eqref{eq:annihilator-full-linearization}分解为
\begin{align}
 \int_\Sigma\langle D_1\Phi(u,J)\xi,\eta\rangle&=0,
 \label{eq:annihilator-map-direction}\\
 \int_\Sigma\langle D_2\Phi(u,J)B,\eta\rangle&=0
 \label{eq:annihilator-structure-direction}
\end{align}
对所有 $\xi\in T_u\mathcal F(A)$ 及 $B\in C^\ell$ 成立。第一式等价于分布意义下
\[
 (D_1\Phi(u,J))^\dagger\eta=0,
\]
其中 $(D_1\Phi(u,J))^\dagger$ 是 $D_1\Phi(u,J)$ 的伴随。注意
$(D_1\Phi(u,J))^\dagger$ 是具有 $C^{\ell-1}$ 系数的一阶形式椭圆微分算子，其主符号为
\[
 \sigma(D_u\bar\partial_{(j,J)})^\dagger:
 T^*\Sigma\setminus o_{T^*\Sigma}\longrightarrow
 \operatorname{Hom}_{\mathbb R}
 \bigl(\pi^*(\Lambda^{(0,1)}(u^*TM))^\dagger,
 \pi^*(u^*TM)^\dagger\bigr),
\]
其中 $\pi:T^*\Sigma\to\Sigma$ 是投影，$(\cdot)^\dagger$ 表示对偶空间。因为
$\sigma(D_u\bar\partial_{(j,J)})(\alpha)$ 是复标量型同构，其伴随亦然。
一般椭圆理论（见命题 10.1.13）于是说明 $\eta$ 可微，并且是方程
\begin{equation}
 (D_1\Phi(u,J))^\dagger\eta=0
 \label{eq:adjoint-linearized-equation}
\end{equation}
的经典解。

下面的唯一延拓引理来自纳赫曼·阿龙沙因（Nachman Aronszajn）对拉普拉斯型椭圆算子的唯一延拓引理 (Aro57)：考虑二阶方程
$D_1\Phi(u,J)(D_1\Phi(u,J))^\dagger\eta=0$ 即可。

\begin{lemma}[唯一延拓引理]\label{lem:unique-continuation-adjoint}
方程\eqref{eq:adjoint-linearized-equation}的任意解，若在 $\Sigma$ 的某个开集上恒为零，则处处恒为零。
\end{lemma}

\begin{exercise}\label{ex:unique-continuation-perturbed-cr}
考虑开域 $D\subset\mathbb C$ 上的扰动柯西--黎曼方程
$\bar\partial\eta+A(z)\eta=0$，其中
$\eta:D\to\mathbb C^n$，$A:D\to\operatorname{End}_{\mathbb R}(\mathbb C^n)$ 光滑，证明上述唯一性引理。事实上，只需假设 $A$ 连续；证明见 (FHS95)。
\end{exercise}

因此，只需证明 $\eta$ 在某个开集上为零。此时在\emph{某处单射}条件和第二个方程发挥作用：
\begin{equation}
 \int_\Sigma\langle D_2\Phi(u,J)B,\eta\rangle=0.
 \label{eq:annihilator-structure-repeat}
\end{equation}
因为假设 $u$ 在某处单射，存在 $z_0\in\Sigma$ 使
\begin{equation}
 du(z_0)\ne0,
 \qquad \#\bigl(u^{-1}(u(z_0))\bigr)=1.
 \label{eq:somewhere-injective-point}
\end{equation}
于是有 $z_0$ 的开邻域 $z_0\in V\subset\Sigma$，使得把
\eqref{eq:somewhere-injective-point}中的 $z_0$ 换成任意 $z\in V$ 后仍成立。

现在考察切空间 $T_J\mathcal J_\omega$。回忆 $J\in\mathcal J_\omega$ 是
$\operatorname{End}(TM)$ 的截面，满足
\[
 J^2=-id,
 \qquad \omega(J\cdot,J\cdot)=\omega(\cdot,\cdot).
\]
元素 $B\in T_J\mathcal J_\omega$ 是满足
\[
 JB+BJ=0,
 \qquad \omega(B\cdot,J\cdot)+\omega(J\cdot,B\cdot)=0
\]
的 $\operatorname{End}(TM)$ 截面。第二式正是说 $B$ 关于度量
$g_J=\omega(\cdot,J\cdot)$ 对称。形式地，考虑纤维丛
\[
 \bigcup_{x\in M}S_\omega(x)=S_\omega\longrightarrow M,
\]
其中
\[
 S_\omega(x)=\{B_x\in\operatorname{End}(T_xM)\mid
 J(x)B_x+B_xJ(x)=0,
 \ \omega_x(B_x,J(x))+\omega_x(J(x),B_x)=0\}.
\]
于是 $B\in T_J\mathcal J_\omega$ 定义 $S_\omega$ 的一个截面。注意
$S_\omega(x)$ 同构于向量空间
\[
 S(\mathbb R^{2n})=\{a\in\operatorname{End}(\mathbb R^{2n})\mid
 aJ_0+J_0a=0, a\text{ 对称}\},
\]
这里 $J_0:\mathbb R^{2n}\to\mathbb R^{2n}$ 是标准复结构，
$g_J=\langle\cdot,\cdot\rangle$ 是标准欧氏内积。由此容易得到
$\dim S(\mathbb R^{2n})=n^2+n$。

\begin{lemma}\label{lem:symmetric-antilinear-orbit}
设 $e\in\mathbb R^{2n}$ 是固定非零向量，则
$S(\mathbb R^{2n})\cdot e=\mathbb R^{2n}$；也就是说，任意
$\mathbb R^{2n}$ 中的向量都可写成 $a\cdot e$，其中
$a\in S(\mathbb R^{2n})\subset\operatorname{End}(\mathbb R^{2n})$。
\end{lemma}

\begin{exercise}\label{ex:symmetric-antilinear-orbit}
证明这个引理。
\end{exercise}

必要时缩小 $V$。我们反设不能在 $V$ 上证明 $\eta\equiv0$，则存在
$z_1\in V$ 使 $\eta(z_1)\ne0$。因为当
$\bar\partial_{(j,J)}u(z_1)=0$ 时，$du(z_1)$ 的秩只能是 $0$ 或 $2$（为什么？），所以若
$du(z_1)\ne0$，它必为单射。非恒定 $J$-全纯映射 $u$ 的临界点集是离散的；因此，在 $V$ 中略微移动给定点后，还可以假设
$du(z_1)\ne0$，从而 $du(z_1)$ 是单射。取单位向量
$v_{z_1}\in T_{z_1}\Sigma$，使 $\eta(z_1)(v_{z_1})\ne0$，则
\[
 du(z_1)(jv_{z_1})\ne0.
\]
由上述引理，可取 $a_x\in S_\omega(x)$（$x=u(z_1)$），使
\[
 a_x\cdot du(z_1)(jv_{z_1})=\eta(z_1)(v_{z_1}).
\]

\begin{figure}[htbp]
 \centering
 \begin{tikzpicture}[scale=.92,>=Stealth]
  \draw[smooth] (-4.2,.25) .. controls (-3.6,.05) and (-3.1,.18) .. (-2.55,.22);
  \fill (-3.55,.12) circle (1.3pt);
  \node[below] at (-3.55,.08) {$z_1$};
  \node[above] at (-2.95,.28) {$z_1\in V\subset\Sigma$};
  \draw[->] (-2.25,.2)--(-1.15,.2);
  \node[below] at (-1.7,.2) {$u$};
  \draw[smooth] (-.75,-.35)
    .. controls (-.05,.05) and (.2,.2) .. (.05,-.12)
    .. controls (-.12,-.42) and (.18,-.55) .. (.65,-.45)
    .. controls (1.2,-.32) and (1.15,.28) .. (1.85,.2)
    .. controls (2.35,.15) and (2.75,.18) .. (3.35,.0);
  \draw (1.35,.18) circle (.48);
  \fill (1.35,.18) circle (1.4pt);
  \node[below right] at (1.35,.18) {$u(z_1)$};
  \node[below] at (1.35,-.3) {$U$};
  \node[right] at (3.35,.0) {$\operatorname{Im}(u)$};
\end{tikzpicture}

 \caption{嵌入点的邻域。}
 \label{fig:neighborhood-embedded-point}
\end{figure}

由于 $u$ 在 $V$ 上是嵌入且 $z_1\in V$，存在 $x=u(z_1)$ 的开邻域
$U\subset M$，使 $u^{-1}(U)=V$。由 $v_{z_1}$ 与 $a_x$ 的选取，有
\[
 \bigl\langle a_x\cdot du(z_1)(jv_{z_1}),
 \eta(z_1)(v)\bigr\rangle
 =|\eta(z_1)(v_{z_1})|^2>0.
\]
由于 $\{v_{z_1},jv_{z_1}\}$ 构成 $T_{z_1}\Sigma$ 的标准正交基，而 $\eta$ 是关于
$(j,J)$ 的 $(0,1)$-形式，由算子内积的定义得到
\begin{align*}
 \langle a_x\cdot du(z_1)(j),\eta(z_1)\rangle
 &:={}&\langle a_x\cdot du(z_1)(j)(v_{z_1}),\eta(z_1)(v_{z_1})\rangle\\
 &&+\langle a_x\cdot du(z_1)(j)(jv_{z_1}),\eta(z_1)(jv_{z_1})\rangle\\
 &=2\langle a_x\cdot du(z_1)(jv_{z_1}),\eta(z_1)(v_{z_1})\rangle\\
 &=2|\eta(z_1)(v_{z_1})|^2>0.
\end{align*}
倒数第二个等号用到了 $du$ 与 $\eta$ 都关于 $(j,J)$ 复线性这一性质。见图 10.2。

取局部截面 $B\in\Gamma(S_\omega)$，使 $B(x)=a_x$ 且
$\operatorname{supp}B\subset U$。若取 $U'\subset U$ 充分小，使
$u(z_1)\in U'$ 且 $u^{-1}(\overline U')\subset V$，则对所有
$z\in u^{-1}(\overline U')$ 有
\[
 \langle B(u(z))du(z)\circ j,\eta(z)\rangle>0.
\]
再取截断函数 $\chi:M\to\mathbb R$，满足
\[
 \chi\equiv1\text{ 于 }U',
 \qquad \operatorname{supp}\chi\subset U.
\]
若把 $\chi$ 取得在 $L^1$ 中充分接近 $U'$ 的特征函数 $\chi_{U'}$，则由
$D_2\Phi(u,J)$ 的公式\eqref{eq:variation-almost-complex-structure}有
\[
 0<\int_\Sigma\left\langle\frac12(\chi B)(u)\cdot du\circ j,\eta\right\rangle
 =\int_\Sigma\langle D_2\Phi(u,J)(\chi B),\eta\rangle.
\]
但把 $\delta J=\chi B$ 代入恒等式
$\int_\Sigma\langle D_2\Phi(u,J)\delta J,\eta\rangle=0$，便得到矛盾。这就证明
$D\Phi(u,J)$ 满射。
\end{proof}

回到定理的证明。既然已经证明 $D\Phi(u,J)$ 满射，隐函数定理立即说明，在给定
$(u,J)\in\widetilde{\mathcal M}^{\mathrm{inj}}(\beta)$ 附近，
$\widetilde{\mathcal M}(\beta)$ 是
$\mathcal F^{k,p}(A)\times\mathcal J_\omega^\ell$ 的（无穷维）光滑子流形。因此定理 10.4.1 的第 (1) 项成立。

现在取定理所给的
$\widetilde{\mathcal M}^{\mathrm{inj}}(\beta)\subset
\mathcal F^{k,p}(\beta)\times\mathcal J_\omega^\ell$ 中 $(u,J)$ 的一个邻域，并考虑投影
\[
 \Pi_\beta:=\left.\pi_2\right|_{\widetilde{\mathcal M}^{\mathrm{inj}}(\beta)}
 \longrightarrow\mathcal J_\omega,
 \qquad \Pi_\beta(u,J)=J.
\]
以 $\iota(g;\beta)$ 表示与
$(j,u)\in\mathcal J(\Sigma)\times\mathcal F_g(M,\omega;\beta)$（其中 $g_\Sigma=g$）相伴的拓扑指标。它与任意
$J\in\mathcal J_\omega$ 下 $D_{(j,u)}\bar\partial_J$ 的解析指标一致，由黎曼--罗赫公式可算得
\[
 \operatorname{Index}D\bar\partial_{(j,J)}(u)=\iota(g;\beta),
\]
其中
\begin{equation}
 \iota(g;\beta)=
 \begin{cases}
  2\bigl(c_1(M,\omega)(\beta)+n(1-g)\bigr),&g\geq2,\\
  2\bigl(c_1(M,\omega)(\beta)+1\bigr),&g=1,\\
  2\bigl(c_1(M,\omega)(\beta)+n\bigr),&g=0.
 \end{cases}
 \label{eq:topological-index-by-genus}
\end{equation}
定理 10.4.1 的第 (2) 项可改述如下。

\begin{proposition}\label{prop:fredholm-projection-universal-moduli}
$\Pi_\beta:\widetilde{\mathcal M}(\beta)\to\mathcal J_\omega$ 是弗雷德霍姆映射，且
\[
 \operatorname{Index}D\Pi_\beta(u,J)
 =\operatorname{Index}D\bar\partial_{(j,J)}(u).
\]
见图 10.3。
\end{proposition}

\begin{figure}[htbp]
 \centering
 \begin{tikzpicture}[scale=.92,>=Stealth]
  \draw[->] (-3,-1.2)--(3.4,-1.2);
  \draw[->] (-2.5,-1.65)--(-2.5,1.65);
  \draw[->] (-2.95,.55)--(-2.95,-.95);
  \node[right] at (3.4,-1.2) {$\mathcal J_\omega$};
  \node[left] at (-2.5,1.55) {$\mathcal F(\beta)$};
  \node[left] at (-2.95,-.2) {$\Pi_\beta$};
  \draw[smooth] (-.95,-.8) .. controls (-.65,-.35) and (.3,-.35) .. (1.35,-.7)
    .. controls (1.75,-.85) and (2.15,-.9) .. (2.65,-.9);
  \draw[smooth] (-.2,.8) .. controls (.65,.85) and (.55,.35) .. (-.15,.15)
    .. controls (-.7,-.05) and (-.5,-.5) .. (.25,-.45)
    .. controls (.9,-.43) and (1.35,-.48) .. (1.95,-.5);
\end{tikzpicture}

 \caption{投影 $\Pi_\beta$。}
 \label{fig:projection-pi-beta}
\end{figure}

结合萨德--斯梅尔定理，立刻得到下面的推论。

\begin{corollary}\label{cor:dense-regular-almost-complex-structures}
存在稠密集 $\mathcal J_{\mathrm{reg}}^\ell\subset\mathcal J_\omega^\ell$，使对
$J\in\mathcal J_{\mathrm{reg}}^\ell$，
$\widetilde{\mathcal M}^{\mathrm{inj}}(J;\beta)$ 是维数为
$\iota(g;\beta)$ 的光滑流形。
\end{corollary}

\begin{proof}[命题 10.4.7 的证明]
注意 $D\Pi_\beta:T_{(u,J)}\widetilde{\mathcal M}(\beta)\to
T_J\mathcal J_\omega$ 只是投影的限制：对
$(\xi,B)\in T_{(u,J)}\widetilde{\mathcal M}(\beta)=\ker D\Phi(u,J)$，有
$D\Pi_\beta(\xi,B)=B$。核 $\ker D\Pi_\beta(u,J)$ 有限维，因为
\begin{align*}
 \ker D\Pi_\beta(u,J)
 &=\ker D\Phi(u,J)\cap\{T_u\mathcal F\times\{0\}\}\\
 &=\ker D_1\Phi(u,J)\\
 &=\ker D\bar\partial_{(j,J)}(u).
\end{align*}
下面证明
$\operatorname{coker}D\Pi_\beta(u,J)\cong
\operatorname{coker}(D\bar\partial_Ju)$。考虑下列图表：
\[
\begin{tikzpicture}[>=Stealth,scale=.77,transform shape]
 \node (zeroL) at (-5,0) {$0$};
 \node (ker) at (-3.5,0) {$\ker D\Phi$};
 \node (sum) at (-.8,0) {$T_u\mathcal F\oplus T_J\mathcal J_\omega$};
 \node (omega) at (3.0,0) {$\Omega_J^{(0,1)}(u^*TM)$};
 \node (zeroR) at (5.7,0) {$0$};
 \node (tf) at (-.8,1.8) {$T_u\mathcal F$};
 \node (tj) at (-.8,-1.8) {$T_J\mathcal J_\omega$};
 \node (cokerp) at (3.0,-2.5) {$\operatorname{coker}D\Pi_\beta$};
 \node (cokerd) at (5.2,-1.35) {$\operatorname{coker}D\bar\partial_{(j,J)}(u)$};
 \draw[->] (zeroL)--(ker);
 \draw[->] (ker)--(sum);
 \draw[->] (sum)--node[above] {$D\Phi$}(omega);
 \draw[->] (omega)--(zeroR);
 \draw[->] (sum)--node[left] {$\pi_1$}(tf);
 \draw[->] (tf)--node[above right] {$D_1\Phi$}(omega);
 \draw[->] (sum)--node[left] {$\pi_2$}(tj);
 \draw[->] (tj)--node[below right] {$D_2\Phi$}(omega);
 \draw[->] (ker)--node[below left] {$D\Pi_\beta$}(tj);
 \draw[->] (tj)--(cokerp);
 \draw[->] (omega)--(cokerd);
 \draw[->,dotted] (cokerp)--(cokerd);
\end{tikzpicture}
\]
其中两个坐标投影分别为 $\pi_1,\pi_2$，而
$D\Phi=D_1\Phi+D_2\Phi$。只要证明存在自然同构
\[
 \operatorname{coker}D\Pi_\beta
 \longrightarrow\operatorname{coker}D\bar\partial_{(j,J)}(u),
\]
也就是
\[
 T_J\mathcal J_\omega^\ell/\operatorname{Im}(D\Pi_\beta)
 \cong
 \Omega_J^{(0,1)}(u^*TM)/
 \operatorname{Im}(D\bar\partial_{(j,J)}(u)),
\]
证明便告完成。

先证良定义。定义
\[
 T_J\mathcal J_\omega^\ell/\operatorname{Im}D\Pi_\beta
 \longrightarrow
 \Omega_J^{(0,1)}(u^*TM)/\operatorname{Im}D\bar\partial_{(j,J)}(u),
 \qquad [B]\longmapsto[D_2\Phi(u,J)B].
\]
回忆 $D_1\Phi(u,J)\xi=D\bar\partial_{(j,J)}(u)\xi$。若
$B\in\operatorname{Im}(D\Pi_\beta)$，即对某个
$\xi\in\Omega^0(u^*TM)$ 有
$(\xi,B)\in T_{(u,J)}\widetilde{\mathcal M}=\ker D\Phi(u,J)$，则
\[
 D_1\Phi(u,J)\xi+D_2\Phi(u,J)B=0,
 \qquad D_2\Phi(u,J)B=-D_1\Phi(u,J)\xi.
\]
右端属于
$\operatorname{Im}D_1\Phi(u,J)=\operatorname{Im}D\bar\partial_{(j,J)}(u)$，故该映射良定义。

再证单射性。若
$D_2\Phi(u,J)(B)\in\operatorname{Im}(D_u\bar\partial_{(j,J)})$，即存在
$\xi\in T_u\mathcal F$ 使
$D_2\Phi(u,J)(B)=D_u\bar\partial_{(j,J)}\xi$，则
\[
 D_2\Phi(u,J)(B)+D_1\Phi(u,J)(-\xi)=0,
\]
所以
$(-\xi,B)\in\ker D\Phi(u,J)=T_{(u,J)}\widetilde{\mathcal M}$。于是
$B\in\operatorname{Im}D\Pi_\beta(u,J)$，即 $[B]=0$。

最后证满射性。再次由哈恩--巴拿赫定理，只须证明满足
\[
 \langle\operatorname{Im}D_1\Phi(u,J),\eta\rangle
 =0=\langle\operatorname{Im}D_2\Phi(u,J),\eta\rangle
\]
的 $\eta$ 必为零。但这等价于
$\langle\operatorname{Im}D\Phi(u,J),\eta\rangle=0$；而
$D\Phi(u,J)$ 已证满射，故 $\eta\equiv0$。
\end{proof}

现在说明如何证明定理 10.4.1 关于 $\mathcal J_\omega$ 的
$C^\infty$ 版本。有两种办法。

\subsection*{弗洛尔的方法}
\stfhSubIndex{I291}{heng^jie^xing}{横截性（transversality）}{fu^luo^er^fang^fa}{弗洛尔方法（Floer’s approach）}

一种办法是沿用弗洛尔在 (Fl88a) 中给出的横截性证明。对每个给定的
$J_0$，取正数序列 $\epsilon=\{\epsilon_k\}_{k=0}^\infty$，并定义可写成
\[
 J=J_0\exp(-J_0B)
\]
的 $J$ 所成的空间 $\mathcal U_\epsilon(J_0)$；其中
$B\in\operatorname{End}(TM)$ 满足：
\begin{enumerate}[label=(\roman*)]
 \item 对所有 $v,w$，$\omega(Bv,w)+\omega(v,Bw)=0$；
 \item $\sum_{k=0}^\infty\epsilon_k|D^kB|<\infty$；
 \item $J$ 驯服 $\omega$，即 $\omega$ 是 $J$-正的。
\end{enumerate}
弗洛尔在 (Fl88a) 中的引理 5.1 证明：若 $\epsilon$ 衰减得足够快，则
$\mathcal U_\epsilon(J_0)$ 成为巴拿赫流形，并且在
$\mathcal J_\omega$ 中 $J_0$ 的某个 $L^2$ 邻域内稠密。直接在
$\mathcal U_\epsilon(J_0)$ 上应用上述横截性证明，得到稠余子集
$\mathcal U_\epsilon(J_0)^{\mathrm{reg}}\subset\mathcal U_\epsilon(J_0)$，使任意
$J\in\mathcal U_\epsilon(J_0)^{\mathrm{reg}}$ 的所有在某处单射 $J$-全纯映射均横截。让
$J_0\in\mathcal J_\omega$ 变化，并对所有 $J_0$ 取
$\mathcal U_\epsilon(J_0)^{\mathrm{reg}}$ 的并，便在
$\mathcal J_\omega$ 中得到一个第二纲子集
$\mathcal J_{\mathrm{Fl}}^{\mathrm{reg}}$，所需横截性在其上成立。

\subsection*{陶布斯的方法}
\stfhSubIndex{I293}{heng^jie^xing}{横截性（transversality）}{tao^bu^si^fang^fa}{陶布斯方法（Taubes’ approach）}

另一种方法最初归于克利福德·亨利·陶布斯（Clifford Henry Taubes）；下面加以说明。麦克达夫与萨拉蒙也在
(MSa94, MSa04) 中阐述过它。

以 $\widetilde{\mathcal M}^{k,\mathrm{inj}}(M)$ 表示由三元组
$((\Sigma,j),u,J)$ 构成的普适模空间，其中 $u$ 在某处单射而 $J\in C^k$。把亏格为
$g$ 的紧曲面之复结构模空间 $\mathcal M_g$ 写成穷竭可数并
\[
 \mathcal M_g=\bigcup_{l=1}^\infty\mathcal U_{g;l},
\]
其中 $\overline{\mathcal U}_{g;l}$ 紧且
$\overline{\mathcal U}_{g;l}\subset\mathcal U_{g;l+1}$。
下面的关键定义使我们能够在 $(G,K;L)$ 的每个有限阶段逐步证明。

\begin{definition}\label{def:taubes-finite-stage-moduli}
对每个给定的 $(G,K;L)\in\mathbb N^3$，选取充分大的
$k=k(G,K)\in\mathbb N$，并考虑子集
\[
 \widetilde{\mathcal M}_{(G,K;L)}^{k,\mathrm{inj}}(M)
 \subset\widetilde{\mathcal M}^{k,\mathrm{inj}}(M),
\]
它由满足下列条件的三元组
$((\Sigma,j),u,J)\in\widetilde{\mathcal M}^{k,\mathrm{inj}}(M)$ 构成：
\begin{enumerate}[label=(\arabic*)]
 \item $g_\Sigma\leq G$ 且 $l\leq L$；
 \item $j\in\overline{\mathcal U}_{g;l}\subset\mathcal M_g$；
 \item $u$ 满足 $\omega([u])\leq K$ 且 $\|du\|_\infty\leq K$；
 \item 存在一个在某处单射点 $z\in\Sigma$，使
 \begin{equation}
  \inf_{z'\ne z}\frac{\operatorname{dist}(u(z),u(z'))}
  {\operatorname{dist}(z,z')}>\frac1K.
  \label{eq:taubes-uniform-injectivity}
 \end{equation}
\end{enumerate}
\end{definition}

最后一项有意义，因为 $z$ 是在某处单射点；由 $u$ 在某处单射的假设，这样的点存在。进一步，每个在某处单射的
$J$-全纯映射都对某个 $(G,K;L)$ 满足 (1)--(4)。以
$\mathcal J_{(G,K;L)}^{k,\mathrm{reg}}\subset\mathcal J_\omega^k$ 表示投影
\[
 \widetilde{\mathcal M}_{(G,K;L)}^{k,\mathrm{inj}}(M)
 \longrightarrow\mathcal J_\omega^k
\]
的正则值。还要选取 $k$，使得当
$\min\{G,K,L\}\to\infty$ 时 $k\to\infty$。
下述命题的证明留作习题。

\begin{proposition}\label{prop:taubes-open-dense-regular-values}
对每个 $(G,K;L)\in\mathbb N^3$，如上选取
$k=k(G,K;L)$。定义
$\mathcal J_{(G,K;L)}^{k,\mathrm{reg}}$ 为投影
$\widetilde{\mathcal M}_{(G,K;L)}^{k,\mathrm{inj}}(M)\to
\mathcal J_\omega^k$ 的正则值，并令
\[
 \mathcal J_{(G,K;L)}^{\mathrm{reg}}
 =\mathcal J_{(G,K;L)}^{k,\mathrm{reg}}\cap\mathcal J_\omega.
\]
则 $\mathcal J_{(G,K;L)}^{\mathrm{reg}}$ 在
$\mathcal J_\omega$ 的 $C^\infty$ 拓扑中既开又稠密。
\end{proposition}

\begin{exercise}\label{ex:taubes-open-dense-regular-values}
证明这个命题。
\end{exercise}

现在定义
\[
 \mathcal J_{Ta}^{\mathrm{reg}}=
 \bigcap_{(G,K)\in\mathbb N^2}\
 \bigcap_{L\in\mathbb N}\mathcal J_{(G,K;L)}^{\mathrm{reg}}.
\]
由于每个
$\mathcal J_{(G,K;L)}^{\mathrm{reg}}\subset\mathcal J_\omega$ 在
$C^\infty$ 拓扑中既开又稠密，此交集属于第二纲，因而证明完成。

\begin{remark}\label{rem:taubes-category-subtlety}
陶布斯方法所用的论证颇为微妙：一方面，应当使用巴拿赫流形
$\mathcal J_\omega^k$，而不是赋予 $C^\infty$ 拓扑的
$\mathcal J_\omega$；另一方面，上述横截性证明必须对 $M$ 的每个同调类和定义域的每个亏格 $g$ 应用。
遗憾的是，不能在 $k\to\infty$ 时对
$\mathcal J_{g,\beta}^{k,\mathrm{reg}}$ 应用归纳论证，因为选出的稠密子集
$\mathcal J_{g,\beta}^{k,\mathrm{reg}}\subset\mathcal J_\omega^k$ 并非开集，因而一般不具有巴拿赫流形结构。而且，对
$\beta,g$ 取交得到的
$\mathcal J_{g,\beta}^{k,\mathrm{reg}}\subset\mathcal J_\omega$ 不一定是
$\mathcal J_\omega$ 中的第二纲子集。因此，不能直接对 $\beta,g$ 取这些集合的交来构造所需稠密子集。不过，借助上面定义的
$\widetilde{\mathcal M}_{(G,K;L)}^{k,\mathrm{inj}}(M)$，上述命题保证
$\mathcal J_{(G,K;L)}^{k,\mathrm{reg}}\cap\mathcal J_\omega$ 在
$\mathcal J_\omega$ 中不但稠密而且开，从而可以取交，得到所需的第二纲子集
$\mathcal J_{Ta}^{\mathrm{reg}}$。
\end{remark}
\stfhSubIndexRangeEnd{I292}{heng^jie^xing}{横截性（transversality）}{ying^she^heng^jie^xing}{映射横截性（mapping transversality）}
